\documentclass[11pt]{article}

\usepackage[margin=1in]{geometry}
\usepackage{amsmath,amssymb,amsfonts,amsthm}
\usepackage{array}
\usepackage{bm}
\usepackage{graphicx}
\usepackage{multirow}
\usepackage{multicol}
\usepackage{nicefrac}
\usepackage{paralist}
\usepackage{enumitem}
\usepackage{verbatim}
\usepackage{thm-restate}
\usepackage[normalem]{ulem}
\usepackage[compact]{titlesec}
\usepackage{tikz}
\usetikzlibrary{arrows}
\usetikzlibrary{arrows.meta}
\usetikzlibrary{backgrounds}
\usetikzlibrary{calc}
\usetikzlibrary{decorations.markings}
\usetikzlibrary{fit}
\usetikzlibrary{patterns}
\usetikzlibrary{positioning}
\usetikzlibrary{shapes}

\usepackage{fancybox}
\newenvironment{fminipage}%
  {\begin{Sbox}\begin{minipage}}%
  {\end{minipage}\end{Sbox}\fbox{\TheSbox}}

\newenvironment{algbox}[0]{%
\vskip 0.2in
\noindent
\begin{fminipage}{6.3in}
}{%
\end{fminipage}
\vskip 0.2in
}

\usepackage{tcolorbox}
\tcbuselibrary{skins,breakable}
\tcbset{enhanced jigsaw}
\usepackage{mdframed}
\usepackage{xcolor}
\usepackage{cite}
\usepackage{url}
\usepackage{xspace}
\usepackage[mathscr]{euscript}
\usepackage{nameref}
\usepackage[
  linktocpage=true,
  pagebackref=true,
  colorlinks=true,
  linkcolor=blue,
  citecolor=blue,
  urlcolor=blue,
  bookmarks,
  bookmarksopen,
  bookmarksnumbered
]{hyperref}
\usepackage[noabbrev,nameinlink]{cleveref}

\newtheorem{theorem}{Theorem}[section]
\newtheorem{example}{Example}[section]
\newtheorem{lemma}{Lemma}[section]
\newtheorem{proposition}[lemma]{Proposition}
\newtheorem{corollary}[lemma]{Corollary}
\newtheorem{claim}[lemma]{Claim}
\newtheorem{fact}[lemma]{Fact}
\newtheorem{conjecture}[lemma]{Conjecture}
\newtheorem{definition}[lemma]{Definition}
\newtheorem{problem}{Problem}
\newtheorem{remark}[lemma]{Remark}
\newtheorem{observation}[lemma]{Observation}

\newtheorem*{theorem*}{Theorem}
\newtheorem*{claim*}{Claim}
\newtheorem*{proposition*}{Proposition}
\newtheorem*{lemma*}{Lemma}
\newtheorem*{problem*}{Problem}

\allowdisplaybreaks
\setlength{\parskip}{3pt}
\setlength{\emergencystretch}{3em}

\newcommand{\BandedPSDEigenpairs}{%
  \textnormal{\textsc{BandedPSD-\hspace{0pt}Eigenpairs}}}
\newcommand{\RandomizedBandedBisect}{%
  \textnormal{\textsc{Randomized\hspace{0pt}Banded\hspace{0pt}Bisect}}}
\newcommand{\ExactSmallBlock}{%
  \textnormal{\textsc{Exact\hspace{0pt}Small\hspace{0pt}Block}}}
\newcommand{\ShiftedInversePower}{%
  \textnormal{\textsc{Shifted\hspace{0pt}Inverse\hspace{0pt}Power}}}

\title{Selected Eigenpairs of Fixed-Bandwidth Positive Definite Matrices in Bit Complexity}
\author{}
\date{August 4, 2026}

\begin{document}

\pagenumbering{roman}

\maketitle

\begin{abstract}
Let $A$ be an $n\times n$ real symmetric positive definite matrix of fixed
half-bandwidth $w$, with nonzero integer entries of $O(\log n)$ bits and
spectrum contained in $[n^{-10},n^{10}]$.  Given $1\leq k\leq n$ and dyadic
$0<\epsilon,\delta\leq1/4$, we give a randomized finite-bit algorithm that
terminates on every random branch and, with probability at least
$1-9\delta/64>1-\delta$, returns approximations to the leading $k$
eigenpairs, equivalently the leading symmetric singular triples.  Each
returned unit vector $q_i$ is represented by a nonzero
finite dyadic vector $y_i$, interpreted as $y_i/\sqrt{y_i^Ty_i}$, and the
nonnegative dyadic eigenvalue estimate $\widehat\sigma_i$ satisfies
\[
 \lVert Aq_i-\widehat\sigma_iq_i\rVert_2
   \leq C_w\epsilon\lVert A\rVert_2,
 \qquad
 |q_i^Tq_j|\leq C_w\epsilon\quad(i\ne j),
 \qquad
 |\widehat\sigma_i-\lambda_i(A)|
   \leq C_w\epsilon\lVert A\rVert_2.
\]
For $Q=[q_1\ \cdots\ q_k]$ the algorithm also guarantees
$\lVert Q^TQ-I\rVert_2\leq C_wk\epsilon$.  If
$L=1+b_\epsilon+b_\delta+\lceil\log_2(n/(\epsilon\delta))\rceil$, where
$b_\epsilon$ and $b_\delta$ are the canonical binary encoding lengths, its
worst-case cost is $O_w(nkL^3)$ bit operations, it uses
$O(nL+kL^2)$ random bits, and its output uses $O(nkL)$ bits.  The
special case in which $\epsilon^{-1}$ and $\delta^{-1}$ are polynomial in
$n$ and their encodings have $O(\log n)$ bits therefore costs
$O_w(nk\log^3 n)$ bit operations.  No input eigengap is assumed.  The
construction uses a finite-bit diagonal perturbation, guarded banded Sturm
bisection, and a multishift bank of shifted-inverse polynomial filters.
Constants denoted by $C_w$ or hidden by $O_w(\cdot)$ depend only on $w$.
\end{abstract}

\noindent\textbf{Note.} This paper was fully AI generated by GPT 5.6 Sol at
ultra reasoning effort using an internal harness, from a fairly detailed prompt
by Richard Peng that outlined the projection and iteration approach and
explicitly requested a bit-complexity analysis rather than a word-complexity
analysis.

\clearpage
\tableofcontents
\clearpage

\pagenumbering{arabic}

\section{Introduction}
\label{sec:introduction}

Computing a selected part of the spectrum of a structured symmetric matrix is
an output-sensitive numerical problem.  An $n\times n$ matrix of fixed
half-bandwidth has only $O_w(n)$ nonzero entries, whereas $k$ explicit
eigenvectors already contain $nk$ coordinates.  At the same time, locating
the corresponding eigenvalues is not enough: the returned vectors should
have small residuals and be mutually nearly orthogonal, even when the input
spectrum is clustered.  In a bit model, the precision used by the numerical
primitives, the randomness, and the finite representation of the output must
all be included in the complexity bound.

For half-bandwidth one, this is the classical selected symmetric
tridiagonal eigenproblem.  Sturm bisection and inverse iteration give the
basic localization-and-recovery paradigm
\cite{barthMartinWilkinson1967Bisection,wilkinson1962InverseIteration}.
MRRR methods target $k$ selected eigenpairs in $O(nk)$ floating-point
arithmetic when the required relatively robust representation properties hold
\cite{parlettDhillon2000RelativelyRobust,dhillonParlett2004MultipleRepresentations,dhillonParlettVomel2006DesignMRRR,willemsLang2013FrameworkMR3}.
Subset calculations and tightly clustered spectra introduce additional
representation and orthogonality issues
\cite{marquesParlettVomel2006Subsets,dhillonParlettVomel2005GluedMatrices,demmelEtAl2008TridiagonalEigensolvers}.
These results treat the half-bandwidth-one special case in floating-point
arithmetic, rather than giving an end-to-end finite-bit bound for the wider
fixed-bandwidth input considered here.  For symmetric band and
fixed-block-width matrices more generally, direct divide-and-conquer and
block-tridiagonal methods compute full or approximate eigensystems in
arithmetic models
\cite{arbenz1992Bandsymmetric,ganstererEtAl2001SymmetricBand,ganstererEtAl2003BlockTridiagonal};
structured spectral divide-and-conquer can instead keep the eigenvector
matrix implicit or compressed
\cite{vogelEtAl2016StructuredEigenvalue,susnjaraKressner2021Banded}.
Those output models do not by themselves give a bit bound for $k$ explicit
vectors.

This question is also a structured analogue of the bidiagonal singular-value
problem posed by Amsel et al.~\cite[Problem~3.8]{amselEtAl2026OpenQuestions},
which asks for $k$ singular triplets of an
arbitrary upper-bidiagonal matrix in $O(nk)$ operations, together with
finite-precision residual and orthogonality guarantees.  MRRR-based and
associated-tridiagonal approaches have been developed for selected
bidiagonal singular triplets
\cite{willemsLangVomel2006BidiagonalSVD,willemsLang2012MR3GK,marquesDemmelVasconcelos2020BidiagonalSVD}.
This is a nearby but different problem: a general bidiagonal input requires
two coupled families of left and right singular vectors.  For a positive
definite symmetric matrix, by contrast, an eigenpair $(\lambda_i,q_i)$
is the symmetric singular triple $(\lambda_i,q_i,q_i)$.  We study this
symmetric problem under fixed bandwidth, polynomial conditioning, and
finite integer input; the result does not solve the arbitrary-bidiagonal
problem.

\paragraph{Our contribution.}
We give and analyze
\hyperref[alg:main]{\BandedPSDEigenpairs}, a randomized finite-bit
algorithm for the leading $k$ eigenpairs in this setting.  All complexity
bounds below count bit operations.  Constants denoted by $C_w$ or hidden by
$O_w(\cdot)$ may depend only on the fixed half-bandwidth $w$.

\begin{theorem}[Top eigenpairs in bit complexity]
\label{thm:main}
Let $A$ be an $n\times n$ real symmetric positive definite matrix with
$A_{ij}=0$ whenever $|i-j|>w$, where $w$ is fixed.  Suppose every nonzero
entry of $A$ is an integer with $O(\log n)$ bits and
\[
 n^{-10}\leq\lambda_n(A)\leq\lambda_1(A)\leq n^{10}.
\]
Given $1\leq k\leq n$ and dyadic $0<\epsilon,\delta\leq1/4$, the randomized
algorithm \hyperref[alg:main]{\BandedPSDEigenpairs} terminates on
every random branch and outputs $k$ triples
$(\widehat\sigma_i,q_i,q_i)$, where $\widehat\sigma_i$ is a nonnegative
dyadic and $q_i$ is a normalized dyadic-vector encoding of a real unit vector.
With probability at least $1-9\delta/64$, and therefore at least
$1-\delta$, simultaneously for $1\leq i\leq k$ and distinct
$i,j\in\{1,\ldots,k\}$,
\begin{align*}
 \lVert Aq_i-\widehat\sigma_iq_i\rVert_2
   &\leq C_w\epsilon\lVert A\rVert_2,\\
 |q_i^Tq_j|&\leq C_w\epsilon\qquad(i\ne j),\\
 |\widehat\sigma_i-\lambda_i(A)|
   &\leq C_w\epsilon\lVert A\rVert_2.
\end{align*}
Moreover, for $Q=[q_1\ \cdots\ q_k]$,
$\lVert Q^TQ-I\rVert_2\leq C_wk\epsilon$.  Writing
\[
 L=1+b_\epsilon+b_\delta+
   \left\lceil\log_2\frac{n}{\epsilon\delta}\right\rceil,
\]
where $b_\epsilon$ and $b_\delta$ are the canonical binary encoding lengths
of the two dyadic inputs,
the algorithm uses $O_w(nkL^3)$ bit operations, uses
$O(nL+kL^2)$ random bits, and produces an $O(nkL)$-bit output.  In
particular, if $\epsilon^{-1}$ and
$\delta^{-1}$ are polynomial in $n$ and their encodings have $O(\log n)$
bits, the cost is $O_w(nk\log^3 n)$ bit operations.
\end{theorem}

The theorem assumes no eigengap for $A$.  In particular, when eigenvalues
are repeated, it does not single out a prescribed basis of the corresponding
eigenspace; the top-$k$ guarantee is expressed through indexed eigenvalue
estimates, residuals, and near orthogonality.  When $\epsilon^{-1}$ and
$\delta^{-1}$ are polynomial in $n$ and their canonical encodings have
$O(\log n)$ bits, the running time is $O_w(nk\log^3 n)$ and the explicit
output occupies $O(nk\log n)$ bits.

On the joint gap, safe-localization, and random-overlap event used in the
proof, which has probability at least $1-9\delta/64$, a small finite-bit
random diagonal perturbation creates quantitative separation.  Guarded
banded Sturm bisection then localizes each leading eigenvalue of the
perturbed matrix.  Around every localization, an even power of a shifted
inverse acts as an almost rank-one spectral filter; the sum of these filters
approximates the leading spectral projector.  Applying the individual
filters to a finite random probe recovers the $k$ vectors without forming and
orthogonalizing a dense range basis.  Finally, certified Rayleigh quotients
and perturbation bounds transfer the residuals and indexed eigenvalue
estimates back to the original matrix.  Fixed bandwidth makes every inertia
computation and shifted solve linear in $n$ up to the explicitly counted
precision factors; explicit guards ensure termination on every random
branch.

\paragraph{Nearby techniques and complexity results.}
Polynomial and rational filtering, FEAST, and spectrum-slicing methods also
construct approximate projectors from matrix functions or shifted solves
\cite{tangPolizzi2014FEAST,liEtAl2016ThickRestartFiltering,xiSaad2016LeastSquaresRationalFilters,liEtAl2019EVSL}.
They typically recover a block invariant subspace and use a Ritz or
orthogonalization stage, which is a different output mechanism from the
individually centered filters used here.  The level-spacing estimate behind
our diagonal perturbation comes from random-operator correlation bounds
\cite{bellissardHislopStolz2007CorrelationEstimates}; these are probabilistic
spectral estimates, not eigensolvers.  Work on stable or bit-complexity
diagonalization for general matrices and dense Hermitian inputs provides a
broader complexity context
\cite{banksEtAl2023PseudospectralShattering,shah2025HermitianDiagonalization,sobczyk2025HermitianEigenproblems},
but concerns full, dense, or values/projector output models rather than the
fixed-bandwidth selected output considered here.

\paragraph{Organization.}
\Cref{sec:preliminaries} collects spectral and perturbation notation,
block-banded inertia facts, and the finite-precision model.
\Cref{sec:overview} explains the separation, localization, filtering,
recovery, and certification steps.  \Cref{sec:proof} constructs the numerical
primitives and the main algorithm and proves its correctness, probability,
and bit-complexity bounds.

\section{Preliminaries}
\label{sec:preliminaries}

\paragraph{Spectral notation.}
All vector spaces are real.  The notation $\lVert\cdot\rVert_2$ denotes the
Euclidean norm for vectors and the induced spectral norm for matrices.  For a
non-singular matrix $B$, write
$\kappa_2(B)=\lVert B\rVert_2\lVert B^{-1}\rVert_2$.  For a
real symmetric $n\times n$ matrix $T$, order the eigenvalues as
$\lambda_1(T)\geq\cdots\geq\lambda_n(T)$ and choose a corresponding
orthonormal eigenbasis $u_1(T),\ldots,u_n(T)$.  We write
\[
 \operatorname{spec}(T)=\{\lambda_1(T),\ldots,\lambda_n(T)\},
 \qquad
 \operatorname{dist}(z,\operatorname{spec}(T))
   =\min_i|z-\lambda_i(T)|,
\]
and let $e_i$ be the $i$th coordinate vector.  For symmetric matrices,
$S\preceq T$ means that $T-S$ is positive semidefinite.  Under the hypotheses
of \Cref{thm:main}, $A$ is positive definite.

For a scalar function $f$ defined on $\operatorname{spec}(T)$, the spectral
calculus gives
\begin{equation}
 \label{eq:spectral-calculus}
 f(T)=\sum_{i=1}^n f(\lambda_i(T))u_i(T)u_i(T)^T,
 \qquad
 \lVert f(T)\rVert_2=\max_i|f(\lambda_i(T))|.
\end{equation}
In particular, for an interval $J\subseteq\mathbb R$ we use
\[
 \mathbf 1_J(T)=\sum_{\lambda_i(T)\in J}u_i(T)u_i(T)^T,
 \qquad
 \lVert(T-zI)^{-1}\rVert_2
   =\operatorname{dist}(z,\operatorname{spec}(T))^{-1}
\]
when $z\notin\operatorname{spec}(T)$.  We use Weyl's standard perturbation
bound
\begin{equation}
 \label{eq:weyl}
 |\lambda_i(T+E)-\lambda_i(T)|\leq\lVert E\rVert_2
 \qquad(1\leq i\leq n)
\end{equation}
for symmetric $E$.  We also use the elementary normalization estimate: if
$u$ is unit, $\alpha\ne0$, and
$\lVert y-\alpha u\rVert_2\leq\eta<|\alpha|$, then
\begin{equation}
 \label{eq:normalization-stability}
 \left\lVert\frac{y}{\lVert y\rVert_2}
   -\operatorname{sgn}(\alpha)u\right\rVert_2
 \leq\frac{2\eta}{|\alpha|}.
\end{equation}

\paragraph{Banded matrices and inertia.}
A matrix has half-bandwidth $w$ if $T_{ij}=0$ whenever $|i-j|>w$.
Set $b=\max\{1,w\}$ and partition the coordinates into consecutive blocks of
size $b$, with a possibly shorter final block.  Then every matrix of
half-bandwidth $w$ is block tridiagonal.  We write $T^{[j]}$ for the leading
principal matrix through block $j$, and use the same superscript for a
conformally blocked matrix.  For $\eta>0$, a real query $x$ is
\emph{$\eta$-safe for $T$} if
\begin{equation}
 \label{eq:safe-query}
 \min_j\operatorname{dist}(x,\operatorname{spec}(T^{[j]}))\geq\eta.
\end{equation}

For a symmetric matrix $P$, let $\iota_+(P)$ be its number of positive
eigenvalues, counted with multiplicity.  Thus
\begin{equation}
 \label{eq:inertia-count}
 \iota_+(T-xI)=|\{i:\lambda_i(T)>x\}|.
\end{equation}
We use Sylvester's law of inertia in its block Schur-complement form.  Write
$B_j$ and $C_j$ for the diagonal and superdiagonal blocks of $T$,
respectively.  If the block $LDL^T$ elimination of $T-xI$ has nonsingular
pivots
\begin{equation}
 \label{eq:exact-block-pivots}
 P_1=B_1-xI,
 \qquad
 P_j=B_j-xI-C_{j-1}^TP_{j-1}^{-1}C_{j-1},
\end{equation}
then
\begin{equation}
 \label{eq:block-inertia}
 \iota_+(T-xI)=\sum_j\iota_+(P_j).
\end{equation}
Safety in \Cref{eq:safe-query} guarantees that these exact pivots are
nonsingular.

For the input matrix $A$, define the scale
\[
 W=\max_{1\leq i\leq n}\sum_{j=1}^{n}|A_{ij}|.
\]
Symmetry and fixed bandwidth give
\begin{equation}
 \label{eq:W-comparison}
 \lVert A\rVert_2\leq W\leq(2w+1)\lVert A\rVert_2.
\end{equation}
The first inequality is the symmetric maximum-row-sum bound, while bandwidth
gives the second.  Moreover, $W\geq1$ and $W$ has $O(\log n)$ bits.

\paragraph{Finite representations and guarded arithmetic.}
A finite dyadic is a number $m2^e$ with integers $m,e$.  A
\emph{normalized dyadic-vector encoding} is a nonzero finite dyadic vector
$y$, interpreted as the exact real unit vector $y/\sqrt{y^Ty}$.  Only $y$
and this interpretation are stored.  Homogeneous expressions are evaluated
without expanding the generally irrational normalization; for example,
\[
 \left(\frac{y}{\sqrt{y^Ty}}\right)^TA
 \left(\frac{y}{\sqrt{y^Ty}}\right)=\frac{y^TAy}{y^Ty}.
\]
All running times count bit operations.  We use schoolbook arithmetic, so an
arithmetic operation on $O(p)$-bit integer data costs $O(p^2)$ bit
operations.  The canonical form of a nonzero dyadic is its reduced
mantissa--exponent pair $m2^e$, with $m$ odd.  For the dyadic inputs
$\epsilon$ and $\delta$, let $b_\epsilon$ and $b_\delta$ denote the total
binary lengths of the signed integers in these canonical forms, and put
\[
 L=1+b_\epsilon+b_\delta+
 \left\lceil\log_2\frac{n}{\epsilon\delta}\right\rceil.
\]

Fix the constant
\[
 K_w=2^{20(w+1)^3}.
\]
The following rounding model is shared by all numerical primitives.

\begin{definition}[Guarded dyadic arithmetic]
\label{def:guarded-dyadic-arithmetic}
For integers $p\geq8$ and $E\geq p$, a guarded $p$-bit dyadic number is zero
or a number $\pm m2^e$, where $m,e\in\mathbb Z$,
$2^{p-1}\leq m<2^p$, and $|e|\leq E$.
Every point operation $+,-,\times,\div,\sqrt{\phantom{x}}$ returns the
nearest such dyadic to the exact result, with ties resolved toward an even
significand.  Put
\[
 u_p=2^{1-p},\qquad \tau_{p,E}=2^{p-1-E}.
\]
Absent exponent overflow, the point-rounding map satisfies
\[
 |\operatorname{rnd}_{p,E}(s)-s|
 \leq u_p|s|+\tau_{p,E}.
\]
A result whose magnitude is below the least positive representable magnitude
may underflow to zero; this is not a failure.  A zero divisor, a negative
square-root argument, or exponent overflow returns $\mathrm{FAIL}$.

For interval certificates, lower and upper endpoints are rounded outward.
They enclose the exact result, and each endpoint differs from it by at most
$2u_p$ times its magnitude plus $\tau_{p,E}$ while the exponent guard is
respected; interval division through zero returns $\mathrm{FAIL}$.  For an
interval $I=[\ell,u]$, write $\inf I=\ell$, $\sup I=u$, and
$\operatorname{width}(I)=u-\ell$.  From a certified interval, we select
$\operatorname{rnd}_{p,E}((\ell+u)/2)$ and use the larger endpoint distance
as a certified radius; for a box this rule is applied coordinatewise, with an
upward-rounded Euclidean norm of the radii.

An exact finite dyadic input longer than $p$ bits may be used in its first
rounded operation when the input-size and exponent-margin conditions of the
routine hold; the result is guarded and obeys the same point-rounding bound.
Block-Sturm computations use
$E_{\mathrm{blk}}=K_wp$, while Givens and triangular-solve computations use
$E_{\mathrm{sol}}=K_wnp$.  On any failure, a routine returns its specified
finite zero output.
\end{definition}

\section{Overview}
\label{sec:overview}

The proof of \Cref{thm:main} analyzes
\hyperref[alg:main]{\BandedPSDEigenpairs}.  On input
$(A,k,\epsilon,\delta)$, this algorithm produces normalized dyadic encodings
of $k$ approximate top eigenvectors and nonnegative dyadic Rayleigh estimates
obtained from certified intervals.  Two issues drive the construction.
First, $A$ need not have any eigengap, whereas isolating individual
eigenvectors with shifted-inverse filters requires quantitative separation.
Second, finite-precision Sturm localization must avoid every
leading-principal spectrum, and all subsequent spectral operations must be
certified at a common precision while retaining a cost essentially linear in
both $n$ and $k$.  The proof resolves these issues by creating a harmless gap,
locating the perturbed eigenvalues with safe inertia queries, isolating them
with a bank of shifted-inverse filters, and certifying the resulting vectors
with banded finite-precision computations.

\paragraph{Creating a usable eigengap.}
We measure the matrix at the row-sum scale
$W=\max_i\sum_j|A_{ij}|$, which satisfies
$\lVert A\rVert_2\leq W\leq(2w+1)\lVert A\rVert_2$.  We then sample a
finite dyadic diagonal matrix $D$ and put
\[
 H=A+D,
 \qquad 0\preceq D\preceq\rho I,
 \qquad \rho\leq\frac{\epsilon W}{128}.
\]
The finite-dimensional Minami estimate first shows that a continuous
diagonal perturbation has pairwise eigenvalue gaps at least $4\Delta$ with
high probability.  Sampling instead on a dyadic grid of spacing
$h\leq\Delta/8$ changes every eigenvalue by less than $h$, so Weyl's bound
preserves a gap of at least $3\Delta$ while using only
$O(n\log(n/(\epsilon\delta)))$ random bits.  This perturbation is introduced
only to split spectral clusters: because $\lVert D\rVert_2\leq\rho$, its
effect fits within the final error budget.  In particular, the proof later
transfers residuals and indexed eigenvalue estimates from $H$ to $A$ rather
than claiming that a chosen eigenbasis of a clustered $A$ is preserved.

\paragraph{Stable spectral localization.}
Fixed bandwidth makes $H-xI$ block tridiagonal with constant-size blocks.
When its leading Schur pivots are nonsingular, the block Sturm recurrence
reduces eigenvalue counting to
\[
 P_1=B_1-xI,
 \qquad
 P_j=B_j-xI-C_{j-1}^TP_{j-1}^{-1}C_{j-1},
 \qquad
 \iota_+(H-xI)=\sum_j\iota_+(P_j)
 =|\{i:\lambda_i(H)>x\}|.
\]
The numerical difficulty is that a query can be separated from
$\operatorname{spec}(H)$ but extremely close to the spectrum of a leading
principal block, making one of these Schur pivots nearly singular.
\hyperref[alg:banded-bisect]{\RandomizedBandedBisect} therefore
chooses each query from a finite dyadic lattice in the middle half of its
current interval.  With high probability all adaptive queries avoid every
leading-principal spectrum.  At such a safe query,
\hyperref[alg:small-block]{\ExactSmallBlock} computes the inertia of
each rounded constant-size pivot exactly and its inverse to guarded
precision; the collection of rounded pivots is the exact recurrence for a
perturbation too small to change the full inertia.  Each query consequently
gives the correct count and shrinks the interval by a constant factor.  The
algorithm first localizes
$\lambda_k(H)$ to set a rank-$k$ threshold, and localizes the other top
eigenvalues because each will center one spectral slice below.

\paragraph{A multishift spectral step.}
On the gap and safe-localization event, write
$|\widetilde\lambda_i-\lambda_i(H)|\leq a$ and set
\[
 z_i=\widetilde\lambda_i+\frac{\Delta}{2},
 \qquad
 F_i=\left[-\frac{\Delta}{2}(H-z_iI)^{-1}\right]^d,
 \qquad
 \gamma=\widetilde\lambda_k-\Delta,
 \qquad
 S_\gamma=\sum_{i=1}^kF_i.
\]
Here $d$ is the least positive even integer with
$4^{-d}\leq\zeta/(128n)$, where $\zeta\leq n^{-10}$, and the localization
accuracy is chosen so that $a\leq\zeta\Delta/(2^{10}d)$.  The threshold
$\gamma$ has exactly $k$ eigenvalues of $H$ above it.  In the eigenbasis of
$H$, the multiplier of $F_i$ at $\lambda_j(H)$ is
$(-\Delta/(2(\lambda_j(H)-z_i)))^d$.  It differs from $1$ by at most
$\zeta/128$ when $j=i$, while the $3\Delta$ gap makes its magnitude at most
$4^{-d}\leq\zeta/(128n)$ when $j\ne i$.  Thus, writing
$u_i=u_i(H)$,
\[
 \lVert F_i-u_iu_i^T\rVert_2\leq\frac{\zeta}{128},
 \qquad
 \left\lVert S_\gamma-
 \mathbf 1_{[\gamma,\infty)}(H)\right\rVert_2
 \leq\zeta\leq n^{-10}.
\]
The sum $S_\gamma$ is therefore the desired approximate step function, while
its individual summands are the nearly rank-one objects used to recover the
eigenvectors.

\paragraph{One-vector recovery and certified inverse powers.}
Independently of $H$, the algorithm samples a finite random integer vector
$r$ and uses its normalized encoding $x=r/\lVert r\rVert_2$ when $r\ne0$,
with $x=e_1$ when $r=0$.  A discrete anti-concentration argument gives
$|u_i(H)^Tx|\geq\beta$ simultaneously for the entire eigenbasis with
probability at least $1-\delta/16$.  For each $i$,
\hyperref[alg:inverse-power]{\ShiftedInversePower} takes
$(H,W,z_i,-\Delta/2,\Delta,\zeta,d,x,p)$,
factors the banded matrix $H-z_iI$ once by Givens QR, and performs $d$
banded triangular solves.  It returns a dyadic vector $y_i$ and a certified
radius $R_i$ satisfying
$\lVert y_i-F_ix\rVert_2\leq R_i\leq\zeta/128$, using
$O_w(ndp^2)$ bit operations.  Its central error
certificate propagates the directed residual radii
$\widehat r_\ell,\widehat s_\ell$ through
\[
 \eta_\ell=\widehat s_\ell+g_+\widehat r_\ell,
 \qquad
 R_\ell=g_+R_{\ell-1}+\eta_\ell.
\]
Here $g_+$ is a directed upper rounding of
$1+\zeta/(256d)$.
This converts local solve residuals into an end-to-end bound for the power
$F_ix$.  Writing $\alpha_i=u_i(H)^Tx$, the overlap event gives
$|\alpha_i|\geq\beta$, while the exact-filter and numerical-application
bounds give
$\lVert y_i-\alpha_i u_i(H)\rVert_2\leq\zeta/64$.  The choice
$\zeta\leq\epsilon\beta/64$ therefore makes the normalized encoding
$q_i=y_i/\sqrt{y_i^Ty_i}$ satisfy
$\lVert q_i-\operatorname{sgn}(\alpha_i)u_i(H)\rVert_2\leq\epsilon/8$.
Recovering the vectors from the individual slices avoids constructing and
orthogonalizing a dense $n\times k$ range basis.

\paragraph{Transfer, certification, and cost.}
Closeness to $u_i(H)$ first gives
$\lVert Hq_i-\lambda_i(H)q_i\rVert_2\leq\epsilon W/2$; since $A=H-D$,
\[
 \lVert Aq_i-\lambda_i(H)q_i\rVert_2
 \leq\lVert Hq_i-\lambda_i(H)q_i\rVert_2+\lVert D\rVert_2
 <\epsilon W.
\]
Directed intervals for the homogeneous Rayleigh quotient
$y_i^TAy_i/(y_i^Ty_i)$ then produce the stored nonnegative dyadic
$\widehat\sigma_i$.  The preceding residual, Weyl's bound, and
the norm comparison in \Cref{eq:W-comparison} give the claimed residual and indexed
eigenvalue errors for $A$.  Because the $u_i(H)$ are orthonormal, the same
vector approximation gives both the pairwise inner-product bound and
$\lVert Q^TQ-I\rVert_2\leq C_wk\epsilon$.

Finally, all scalar parameters have $O(L)$-bit encodings, the common
precision is $p=O_w(L)$, and both the bisection depth and filter degree are
$O(L)$.  A fixed-bandwidth inertia query or shifted solve uses
$O_w(np^2)$ bit operations, so the $k$ localizations and $k$ filter
applications are bounded by $O_w(nkL^3)$ bit operations.  Perturbation,
overlap, and query sampling use $O(nL+kL^2)$ random bits, and the normalized vector
records use $O(nkL)$ output bits.  The gap, safe-query, and overlap events
hold jointly with probability at least $1-9\delta/64$; on this event every
claimed guarantee holds, while explicit guards and finite fallback records
make the algorithm terminate on every random branch.

\section{Proof of the Main Theorem}
\label{sec:proof}

The case $n=1$ is immediate, so until the final algorithm we assume $n\geq2$.
We now fix the dyadic scales used throughout the proof: $\rho$ and $\Delta$
control the perturbation and its induced eigengap, $\beta$ and $\zeta$
control random overlap and filter accuracy, and $a$ and $\nu$ control
localization accuracy and query safety.
Let $\rho$ be the largest power of two not exceeding $\epsilon W/128$, and
let $\Delta$ be the largest power of two not exceeding
$\delta\rho^2/(2^{16}Wn^2)$.  Let $\beta$ be the largest power of two not
exceeding $\delta/(32n^2)$, and let $\zeta$ be the largest power of two not
exceeding $\min\{n^{-10},\epsilon\beta/64\}$.  Thus
\begin{align}
 \label{eq:rho-range}
 \frac{\epsilon W}{256}<\rho&\leq\frac{\epsilon W}{128},\\
 \label{eq:gap-parameters}
 \frac{\delta\rho^2}{2^{17}Wn^2}<\Delta
  &\leq\frac{\delta\rho^2}{2^{16}Wn^2},
 &
 \frac{\delta}{64n^2}<\beta&\leq\frac{\delta}{32n^2},\\
 \label{eq:zeta-range}
 \frac12\min\left\{n^{-10},\frac{\epsilon\beta}{64}\right\}
  <\zeta&\leq
 \min\left\{n^{-10},\frac{\epsilon\beta}{64}\right\}.
\end{align}
Let $d$ be the least positive even integer such that
\begin{equation}
 \label{eq:filter-degree}
 4^{-d}\leq\frac{\zeta}{128n},
\end{equation}
and let $a$ be the largest power of two not exceeding
$\zeta\Delta/(2^{10}d)$.  Hence
\begin{equation}
 \label{eq:eigenvalue-accuracy}
 \frac{\zeta\Delta}{2^{11}d}<a
  \leq\frac{\zeta\Delta}{2^{10}d}.
\end{equation}
For the randomized inertia bisections define
\begin{equation}
 \label{eq:bisection-parameters}
 m_{\mathrm{bis}}=
  \left\lceil\log_{4/3}\frac{2W}{a}\right\rceil+2,
 \qquad
 N_{\mathrm{qry}}=(k+1)m_{\mathrm{bis}}.
\end{equation}
Finally, let $\nu$ be the largest power of two not exceeding
$\min\{1,\delta a/(2^{20}n^2N_{\mathrm{qry}})\}$.  Then
\begin{equation}
 \label{eq:query-safety}
 \frac12\min\left\{1,\frac{\delta a}{2^{20}n^2N_{\mathrm{qry}}}\right\}
 <\nu\leq
 \min\left\{1,\frac{\delta a}{2^{20}n^2N_{\mathrm{qry}}}\right\}.
\end{equation}
These definitions give
\begin{align}
 \label{eq:parameter-bits}
 &\log_2(W/\Delta)+\log_2(1/\zeta)+\log_2(W/a)
   +\log_2(W/\nu)
   =O\!\left(\log\frac{n}{\epsilon\delta}\right),\\
 &d+m_{\mathrm{bis}}
   =O\!\left(\log\frac{n}{\epsilon\delta}\right).
\end{align}

We first create the gap used by the multishift construction.  The only
probabilistic input is the following finite-dimensional two-level estimate.

\begin{lemma}[Finite-dimensional Minami estimate]
\label{lem:minami-specialized}
Let $T$ be a fixed real symmetric $n\times n$ matrix, let
$\xi_1,\ldots,\xi_n$ be independent real random variables with a common
density bounded by $K$, and put
$T_\xi=T+\operatorname{diag}(\xi_1,\ldots,\xi_n)$.  For every interval
$J\subseteq\mathbb R$,
\[
 \Pr\!\left\{\operatorname{tr}\mathbf 1_J(T_\xi)\geq2\right\}
 \leq\frac{\pi^2}{2}(Kn|J|)^2.
\]
\end{lemma}

\begin{proof}
This is the $r=2$ finite-volume specialization of the correlation estimate
of Bellissard, Hislop, and Stolz~\cite[Theorem~4]{bellissardHislopStolz2007CorrelationEstimates}.
To obtain the displayed finite-dimensional form, embed $T$ as the finite
restriction of the bounded background operator and apply their estimate on
the coordinate interval $\{1,\ldots,n\}$.  Their bound is
$\Pr\{\operatorname{tr}\mathbf 1_J(T_\xi)\geq r\}
\leq\pi^r(Kn|J|)^r/r!$; substituting $r=2$ gives the displayed inequality.
\end{proof}

The continuous perturbation supplied by the preceding estimate must be
rounded without losing its gap.  The next lemma carries out that finite-bit
step and bounds both its randomness and its norm.

\begin{lemma}[Finite-bit diagonal gap perturbation]
\label{lem:random-gap}
There is a distribution on diagonal dyadic matrices $D$, samplable with
$O(n\log(n/(\epsilon\delta)))$ random bits, such that
$0\preceq D\preceq\rho I$ always and, with probability at least
$1-\delta/16$, the eigenvalues of $H=A+D$ satisfy
\[
 |\lambda_i(H)-\lambda_j(H)|\geq3\Delta
 \qquad(i\ne j).
\]
\end{lemma}

\begin{proof}
Let $\xi_1,\ldots,\xi_n$ be independent uniform random variables on
$[0,\rho]$, and set
$H_c=A+\operatorname{diag}(\xi_1,\ldots,\xi_n)$.  Its spectrum lies in
$[0,W+\rho]\subseteq[0,2W]$.  If two eigenvalues are at distance less than
$4\Delta$, they lie together in one interval of the overlapping cover
\[
 [4j\Delta,4(j+2)\Delta],
 \qquad
 0\leq j\leq\left\lceil\frac{2W}{4\Delta}\right\rceil.
\]
There are at most $W/\Delta+2$ such intervals, each of length $8\Delta$.
Applying \Cref{lem:minami-specialized} with $K=1/\rho$ and taking a union
bound gives
\begin{align*}
 \Pr\!\left\{\min_{i\ne j}|\lambda_i(H_c)-\lambda_j(H_c)|<4\Delta\right\}
 &\leq\left(\frac W\Delta+2\right)
   \frac{\pi^2}{2}\left(\frac{8n\Delta}{\rho}\right)^2\\
 &\leq64\pi^2\frac{Wn^2\Delta}{\rho^2}
 \leq\frac{\pi^2\delta}{2^{10}}<\frac\delta{16},
\end{align*}
where the parameter bounds give $W/\Delta>2$, so
$W/\Delta+2\leq2W/\Delta$, and
\Cref{eq:gap-parameters} is used in the last line.

Choose the least $s$ such that $h=\rho2^{-s}\leq\Delta/8$.  Sample
independent uniform $J_i\in\{0,\ldots,2^s-1\}$ and set $D_{ii}=hJ_i$.
Couple this sample to the continuous one by writing
$\xi_i=h(J_i+U_i)$ with independent $U_i$ uniform on $[0,1)$.  Then
$\lVert D-\operatorname{diag}(\xi_i)\rVert_2<h$.  By Weyl's inequality in
\Cref{eq:weyl}, each eigenvalue moves by less than $h$, so a $4\Delta$ gap of
$H_c$ becomes a gap larger than $4\Delta-2h>3\Delta$ for $H=A+D$.  Finally,
$s=O(\log(n/(\epsilon\delta)))$ by
\Cref{eq:rho-range,eq:gap-parameters}.
\end{proof}

We now specify the finite-precision primitives using the block partition and
guarded arithmetic fixed in the preliminaries.

For later reference, the common sufficient precision condition is
\begin{equation}
 \label{eq:primitive-precision}
 u_p\leq\min\left\{
  \frac{\nu^2}{2^{20}K_wW^2},
  \frac{\zeta\Delta}{2^{20}K_wn^2dW},
  \frac{1}{2^{20}K_wn^2}
 \right\}.
\end{equation}
The first primitive resolves constant-size singular-pivot ambiguities in the
block recurrence by exact rational arithmetic.

\begin{algbox}
\textbf{\uline{\ExactSmallBlock}}: Exact inversion and inertia of a
constant-size symmetric block.
\phantomsection\label{alg:small-block}

\uline{Input}: A symmetric dyadic matrix $P$ of dimension at most $b$ whose
entries are guarded $p$-bit numbers, followed by an integer precision
$p\geq8$; the exponent guard is $E_{\mathrm{blk}}=K_wp$.

\uline{Output}: Either $(\mathrm{OK},X,\iota_+(P))$, where $\iota_+(P)$ is
the exact number of positive eigenvalues of $P$ and $X$ approximates
$P^{-1}$, or $(\mathrm{FAIL},0,0)$.

\uline{Guarantee}: An $\mathrm{OK}$ result has the exact inertia and satisfies
\[
 \lVert X-P^{-1}\rVert_2\leq
 K_w\bigl(u_p\lVert P^{-1}\rVert_2+\tau_{p,E_{\mathrm{blk}}}\bigr).
\]

\begin{enumerate}
 \item Interpret the entries of $P$ as exact rationals.  Compute $\det(P)$
 and $\operatorname{adj}(P)$ from their finite permutation formulas.  If
 $\det(P)=0$, or an exact numerator or denominator exceeds $K_w^2p$ bits,
 return $(\mathrm{FAIL},0,0)$.  Otherwise form
 $P^{-1}=\operatorname{adj}(P)/\det(P)$ exactly and round its entries to
 obtain $X$; an exponent-guard violation returns failure.
 \item Compute inertia by exact symmetric elimination.  At each stage,
 symmetrically permute the least-indexed nonzero diagonal entry into the first
 position and use it as a $1$-by-$1$ pivot.  If all diagonal entries vanish,
 use the lexicographically first nonzero off-diagonal entry $c$ in the pivot
 $\left(\begin{smallmatrix}0&c\\c&0\end{smallmatrix}\right)$.  Add the
 positive inertia of the pivot and recurse on its exact rational Schur
 complement.  Return failure if the bit guard is exceeded; otherwise return
 the accumulated count together with $X$.
\end{enumerate}
\end{algbox}

\begin{lemma}[Exact constant-size block kernel]
\label{lem:small-block-kernel}
\hyperref[alg:small-block]{\ExactSmallBlock} is total, never returns
status $\mathrm{OK}$ with an incorrect positive inertia, satisfies its stated inverse-error bound
whenever it returns $\mathrm{OK}$, and uses $O_w(p^2)$ bit operations.
For the safe block-Sturm pivots below, its status is $\mathrm{OK}$.
\end{lemma}

\begin{proof}
A nonsingular real symmetric matrix always admits the selected pivot.  If no
diagonal entry is nonzero, some off-diagonal entry $c$ is nonzero and the
displayed $2$-by-$2$ pivot has determinant $-c^2$.  A nonsingular principal
pivot in a nonsingular matrix has a nonsingular Schur complement.  Repeated
congruence and inertia additivity therefore prove that the returned inertia
is exact.  The determinant--adjugate identity proves the inverse formula, and
the rounding bound in \Cref{def:guarded-dyadic-arithmetic}, in the fixed
dimension $b$, gives the stated inverse error, including underflow.  The
$K_w^2p$-bit cutoff makes every branch finite.  Only $O_w(1)$ rational
operations are performed, so schoolbook arithmetic on the permitted operands
costs $O_w(p^2)$ bit operations.  Success on safe pivots follows from the
magnitude estimates in \Cref{lem:banded-bisection}.
\end{proof}

The next primitive uses these exact block inertias.  Its query is chosen from
a finite random lattice in the middle half of the current interval, so with
high probability it remains separated from every leading-principal spectrum.
For an interval $I=[\ell,u]$ of length $h_I>2a$, define
\begin{align}
 j^-(I)&=\left\lceil\frac{8(\ell+h_I/4)}{\nu}\right\rceil,
 &
 j^+(I)&=\left\lfloor\frac{8(\ell+3h_I/4)}{\nu}\right\rfloor,\\
 R(I)&=j^+(I)-j^-(I)+1,
 &
 \mathcal X(I)&=
 \left\{\frac\nu8(j^-(I)+J):
 0\leq J<2^{\lfloor\log_2R(I)\rfloor}\right\}.
 \label{eq:bisection-lattice}
\end{align}
A query is safe for the bisection when it is $2\nu$-safe for $T$ in the
sense of \Cref{eq:safe-query}.  If $B_j$ and $C_j$ are the diagonal and
adjacent off-diagonal blocks of $T$, respectively, the point Schur-pivot
recurrence used at query $x$, the rounded counterpart of
\Cref{eq:exact-block-pivots}, is
\begin{equation}
 \label{eq:rounded-block-pivots}
 \widehat P_1=\operatorname{rnd}_{p,E_{\mathrm{blk}}}(B_1-xI),
 \qquad
 \widehat P_j=\operatorname{rnd}_{p,E_{\mathrm{blk}}}
 (B_j-xI-C_{j-1}^TX_{j-1}C_{j-1}),
\end{equation}
where $X_{j-1}$ is the certified rounded inverse returned for the preceding
pivot.  Call a tuple
$(T,i,\allowbreak W,a,\allowbreak \nu,m_{\mathrm{bis}},\allowbreak p)$
\emph{admissible} if $p$ and
$m_{\mathrm{bis}}$ are integers with $p\geq8$, $W,a,\nu$ are finite dyadics,
$W\geq1$,
$T$ is real symmetric of half-bandwidth $w$ with spectrum in $[0,2W]$,
$i\in\{1,\ldots,n\}$,
$0<\nu\leq\min\{a,1\}\leq W$, and
\[
 m_{\mathrm{bis}}\geq
 \left\lceil\log_{4/3}\frac{W}{a}\right\rceil+1,
\]
with every dyadic scalar and nonzero entry of $T$ using at most $K_wp$ bits
and exponent magnitude at most $E_{\mathrm{blk}}/16$, and every integer input
using at most $K_wp$ bits.

\begin{algbox}
\textbf{\uline{\RandomizedBandedBisect}}: Guarded Sturm bisection for
one ordered eigenvalue.
\phantomsection\label{alg:banded-bisect}

\uline{Input}: An admissible tuple
$(T,i,W,a,\nu,m_{\mathrm{bis}},p)$ as defined above.

\uline{Output}: A pair $(\mathrm{status},\widetilde\lambda_i)$ with
$\mathrm{status}\in\{\mathrm{OK},\mathrm{FAIL}\}$; the failure output is
$(\mathrm{FAIL},0)$.

\uline{Guarantee}: Under \Cref{eq:primitive-precision}, if every query is
$2\nu$-safe, then the status is $\mathrm{OK}$ and
$|\widetilde\lambda_i-\lambda_i(T)|\leq a$.

\begin{enumerate}
 \item Initialize $[\ell_0,u_0]=[0,2W]$.
 \item For $r=0,\ldots,m_{\mathrm{bis}}-1$, return the midpoint with status
 $\mathrm{OK}$ if $u_r-\ell_r\leq2a$; otherwise sample $x_r$ uniformly from
 $\mathcal X([\ell_r,u_r])$.  In block order form
 \Cref{eq:rounded-block-pivots} with $x=x_r$ and set
 \[
  (s_{r,j},X_j,\iota_{r,j})
   =\ExactSmallBlock(\widehat P_j,p).
 \]
 If some $s_{r,j}=\mathrm{FAIL}$, return $(\mathrm{FAIL},0)$.  Set
 $N_r=\sum_j\iota_{r,j}$ and update to $[x_r,u_r]$ when $N_r\geq i$, or to
 $[\ell_r,x_r]$ otherwise; name the selected interval
 $[\ell_{r+1},u_{r+1}]$.
 \item After the last update, return the midpoint with status $\mathrm{OK}$
 if the width is at most $2a$, and return $(\mathrm{FAIL},0)$ otherwise.
\end{enumerate}
\end{algbox}

\begin{lemma}[Stable randomized banded Sturm bisection]
\label{lem:banded-bisection}
Under \Cref{eq:primitive-precision},
\hyperref[alg:banded-bisect]{\RandomizedBandedBisect} has its stated
output guarantee and uses
$O_w(nm_{\mathrm{bis}}p^2)$ bit operations and
$O(m_{\mathrm{bis}}\log(W/\nu))$ random bits on every branch.  With the
global parameters in \Cref{eq:gap-parameters,eq:query-safety}, all
at most $N_{\mathrm{qry}}$ adaptive queries are simultaneously safe with
probability at least $1-\delta/64$.
\end{lemma}

\begin{proof}
Fix one invocation and condition on its adaptive history.  If the current
interval has length $h_r>2a$, the full middle-half lattice contains at least
$4h_r/\nu-1$ points.  Its retained power-of-two prefix has at least
$2h_r/\nu-1/2\geq7h_r/(4\nu)$ points.  At most $33$ lattice points can be
within $2\nu$ of one fixed real number.  Thus at most a $64\nu/a$ fraction
of the retained points is unsafe for it.  The sum of the dimensions of all
leading block matrices is at most $n^2$, so the conditional probability that
one query lies within $2\nu$ of any eigenvalue of any $T^{[j]}$ is at most
$64n^2\nu/a$.  The adaptive union bound and \Cref{eq:query-safety} give
\[
 \Pr\{\text{some one of the at most $N_{\mathrm{qry}}$ queries is unsafe}\}
 \leq\frac{64n^2N_{\mathrm{qry}}\nu}{a}<\frac\delta{64}.
\]

Now fix a safe query $x$.  For the point blocks generated by the routine,
define exact symmetric residual blocks
\[
 E_1=\widehat P_1-(B_1-xI),
 \qquad
 E_j=\widehat P_j-
 (B_j-xI-C_{j-1}^T\widehat P_{j-1}^{-1}C_{j-1}).
\]
They include inverse and product rounding.  Inductively, the
$\widehat P_j$ are exactly the block $LDL^T$ pivots of $T+E-xI$, with $E$
block diagonal.  The first residual satisfies
\[
 \lVert E_1\rVert_2\leq
 K_w(u_pW+\tau_{p,E_{\mathrm{blk}}})<\nu/2^{18}.
\]
Safety for $T^{[1]}$ makes $\widehat P_1$ nonsingular with inverse norm below
$\nu^{-1}$ and ensures that its constant-size kernel succeeds.  Suppose the
claim holds through block $j-1$.  Safety and Weyl's inequality put
$T^{[j-1]}+E^{[j-1]}-xI$ at distance at least $31\nu/16$ from singularity.
Its lower-right inverse block is $\widehat P_{j-1}^{-1}$, and therefore
$\lVert\widehat P_{j-1}^{-1}\rVert_2<\nu^{-1}$.  The inverse certificate and
the rounded block product give
\[
 \lVert E_j\rVert_2
 \leq K_w\left(u_p\frac{W^2}{\nu}
   +\tau_{p,E_{\mathrm{blk}}}W^2\right)
   +K_w\tau_{p,E_{\mathrm{blk}}}
 <\frac\nu{2^{18}}.
\]
Safety at the next leading principal matrix now gives
\[
 \lVert\widehat P_j^{-1}\rVert_2<\nu^{-1},
 \qquad
 \lVert\widehat P_j\rVert_2\leq K_w(W+W^2/\nu).
\]
These bounds, \Cref{eq:primitive-precision}, and fixed block dimension show
that all exact rational intermediates have $O_w(p)$ bits, so the block kernel
succeeds and the induction closes.  Consequently
$\lVert E\rVert_2<\nu/16$.

\Cref{eq:inertia-count,eq:block-inertia} say that the
computed count is the number of eigenvalues of $T+E$ above $x$.  Since
$T-xI$ is at distance at least $2\nu$ from singularity, Weyl's inequality
\Cref{eq:weyl} shows that this is also the exact count for $T$.  Every update
therefore retains $\lambda_i(T)$.  A query in the middle half reduces the
interval length by at least a factor $3/4$, and the assumed iteration bound
gives final width at most $2a$.  This proves the output guarantee.

The block recurrence performs $O_w(n)$ guarded scalar operations per query,
each costing $O(p^2)$ bit operations under schoolbook integer arithmetic.
The candidate lattice uses $O(\log(W/\nu))$ random bits per query.  All loops
and failure branches have the stated worst-case bounds.
\end{proof}

We next certify repeated application of a shifted inverse.  For an input
matrix $T$, shift $z$, and scalar $c$, define
\begin{equation}
 \label{eq:shifted-operator}
 B=T-zI,
 \qquad
 G=cB^{-1}.
\end{equation}
Given a point solution $v_\ell$ to $Bv=y_{\ell-1}$, first let $y_\ell$ be
the point rounding of $cv_\ell$.  Then let directed banded matrix-vector
evaluation produce coordinate intervals
$J_{\ell,j}\ni(Bv_\ell-y_{\ell-1})_j$ and
$K_{\ell,j}\ni(y_\ell-cv_\ell)_j$.  With $\operatorname{ru}_p$ denoting
upward rounding under $E_{\mathrm{sol}}$, define
\begin{align}
 \widehat r_\ell&=\operatorname{ru}_p\!\left(
  \sqrt{\sum_{j=1}^n
  \max\{|\inf J_{\ell,j}|,|\sup J_{\ell,j}|\}^2}\right),\\
 \widehat s_\ell&=\operatorname{ru}_p\!\left(
  \sqrt{\sum_{j=1}^n
  \max\{|\inf K_{\ell,j}|,|\sup K_{\ell,j}|\}^2}\right).
 \label{eq:directed-residual-radii}
\end{align}
These quantities certify
\begin{equation}
 \label{eq:solve-certificates}
 \widehat r_\ell\geq\lVert Bv_\ell-y_{\ell-1}\rVert_2,
 \qquad
 \widehat s_\ell\geq\lVert y_\ell-cv_\ell\rVert_2.
\end{equation}
Let $g_+$ be the directed upper rounding of $1+\zeta/(256d)$ and define
\begin{equation}
 \label{eq:radius-recurrence}
 \eta_\ell=\widehat s_\ell+g_+\widehat r_\ell,
 \qquad
 R_\ell=g_+R_{\ell-1}+\eta_\ell.
\end{equation}
These objects turn local residuals into an end-to-end error certificate for
$G^d$.  Call
$(T,W,z,c,\allowbreak\Delta,\zeta,d,\allowbreak x,p)$ an \emph{admissible
shifted-inverse input} if $p$ and $d$ are integers with $p\geq8$ and
$d\geq1$,
$W,z,c,\Delta,\zeta$ are finite dyadics with $W\geq1$, $T$ is real symmetric of
half-bandwidth $w$ with $\lVert T\rVert_2\leq2W$, $x$ is a normalized
dyadic-vector encoding, every dyadic scalar, nonzero entry of $T$, and
coordinate of the representative for $x$ has at most $K_wp$ bits and
exponent in $[-E_{\mathrm{sol}}/16,E_{\mathrm{sol}}/16]$, the integer $d$
has at most $K_wp$ bits, and
\[
 0<\Delta\leq W,\quad 0<\zeta\leq1,\quad
 |c|\leq\Delta,\quad |z|\leq3W,\quad
 \operatorname{dist}(z,\operatorname{spec}(T))\geq\Delta/3,
 \quad
 \lVert G\rVert_2\leq1+\frac{\zeta}{256d}.
\]

\begin{algbox}
\textbf{\uline{\ShiftedInversePower}}: Certified application of an
integer power of a shifted inverse.
\phantomsection\label{alg:inverse-power}

\uline{Input}: An admissible tuple
$(T,W,z,c,\Delta,\zeta,d,x,p)$ as defined above.

\uline{Output}: A triple $(\mathrm{status},y,R)$, where $y$ is a dyadic vector and
$R$ is a nonnegative dyadic error radius; the failure output is
$(\mathrm{FAIL},0,0)$.

\uline{Guarantee}: Under \Cref{eq:primitive-precision}, the status is
$\mathrm{OK}$ and
$\lVert y-G^dx\rVert_2\leq R\leq\zeta/128$.

\begin{enumerate}
 \item Let $r^{\mathrm{enc}}$ be the nonzero dyadic representative of $x$.
 Rescale it by a common exact power of two so that its largest coordinate has
 magnitude in $[1/2,1)$.  Using directed arithmetic, enclose
 $r^{\mathrm{enc}}/\sqrt{(r^{\mathrm{enc}})^Tr^{\mathrm{enc}}}$ in a
 coordinate box, select its midpoint $y_0$, and let $R_0$ be the certified
 Euclidean radius.  If $R_0>\zeta/2048$, return failure.
 \item Form $B$ from \Cref{eq:shifted-operator}.  Compute a point-valued
 banded Givens QR factorization of $B$: after a common exact power-of-two
 rescaling, round each hypotenuse, cosine, and sine to nearest, apply the
 point rotation, using the identity rotation when both entries are zero, and
 set the targeted subdiagonal entry to zero.  Store the
 $O_w(n)$ rotations and the resulting upper-banded matrix $\widehat R$.
 \item For $\ell=1,\ldots,d$, apply the stored rotations to $y_{\ell-1}$ and
 solve the upper-banded system with $\widehat R$ by point back substitution,
 obtaining $v_\ell$.  Compute $\widehat r_\ell,\widehat s_\ell,y_\ell$ from
 \Cref{eq:directed-residual-radii,eq:solve-certificates}, and
 update $\eta_\ell,R_\ell$ by
 \Cref{eq:radius-recurrence}.  Return failure if
 $\eta_\ell>\zeta/(4096d)$, $R_\ell>\zeta/128$, a divisor is zero, or an
 exponent guard is violated.
 \item Return $(\mathrm{OK},y_d,R_d)$.
\end{enumerate}
\end{algbox}

\begin{lemma}[Certified shifted-inverse powering]
\label{lem:shifted-inverse-power}
Under \Cref{eq:primitive-precision},
\hyperref[alg:inverse-power]{\ShiftedInversePower} has its stated
output guarantee and uses $O_w(ndp^2)$ bit operations on every branch.
Successful factorization and solve quantities have exponent magnitude
$O_w(np)$.
\end{lemma}

\begin{proof}
The promises on $B$ imply
\[
 \lVert B\rVert_2\leq5W,
 \qquad
 \lVert B^{-1}\rVert_2\leq3/\Delta,
 \qquad
 \kappa_2(B)\leq15W/\Delta.
\]
We first bound the point-valued QR solve.  One rounded hypotenuse and its two
quotients perturb a $2$-by-$2$ rotation by at most $K_wu_p$ in norm.
Applying it to the $O_w(1)$ live entries in two banded rows introduces a
residual at most $K_wu_p$ times their norm plus
$K_w\tau_{p,E_{\mathrm{sol}}}$.  Telescoping over $O_w(n)$ rotations and
normalizing each computed cosine--sine pair to an exact rotation gives
\[
 \widehat R=\widehat Q^T(B+\Delta B),
 \qquad
 \widehat g=\widehat Q^T(g+\delta g),
\]
where $\widehat Q$ is exactly orthogonal and
\begin{align*}
 \lVert\Delta B\rVert_2
  &\leq K_wn^2(u_pW+\tau_{p,E_{\mathrm{sol}}}),\\
 \lVert\delta g\rVert_2
  &\leq K_wn^2(u_p\lVert g\rVert_2+
        \tau_{p,E_{\mathrm{sol}}}).
\end{align*}

Each back-substitution row has $O_w(1)$ live entries.  Assigning relative
rounding to the corresponding coefficient or right-hand-side entry, and
charging an underflowed quotient after multiplication by its diagonal,
gives
\[
 (\widehat R+\Delta R)\widehat v=\widehat g+\delta h,
\]
with
\begin{align*}
 \lVert\Delta R\rVert_2&\leq K_wn^2u_pW,\\
 \lVert\delta h\rVert_2&\leq
 K_wn^2\bigl(u_p\lVert g\rVert_2+
 W\tau_{p,E_{\mathrm{sol}}}\bigr).
\end{align*}
The input exponent margin implies
\[
 W\tau_{p,E_{\mathrm{sol}}}\leq u_p.
\]
Indeed, $\log_2W\leq K_wp+E_{\mathrm{sol}}/16
\leq9E_{\mathrm{sol}}/16$ for $n\geq2$, and hence
\[
 \log_2\frac{W\tau_{p,E_{\mathrm{sol}}}}{u_p}
 \leq2p-2-\frac{7E_{\mathrm{sol}}}{16}<0.
\]
Multiplication by $\widehat Q$ yields the local backward-error identity
\begin{equation}
 \label{eq:qr-local-backward}
 (B+F)\widehat v=g+f,
 \qquad
 \lVert F\rVert_2\leq4K_wn^2u_pW,
 \qquad
 \lVert f\rVert_2\leq
 2K_wn^2\bigl(u_p\lVert g\rVert_2+
 W\tau_{p,E_{\mathrm{sol}}}\bigr).
\end{equation}

By \Cref{eq:primitive-precision},
$\lVert B^{-1}F\rVert_2<2^{-12}$.  Whenever $\lVert g\rVert_2\leq2$,
\Cref{eq:qr-local-backward} implies
\[
 \lVert\widehat v\rVert_2\leq\frac{12}{\Delta},
 \qquad
 \lVert B\widehat v-g\rVert_2
 \leq\frac{54\zeta}{2^{20}d}
 <\frac{\zeta}{2^{14}d}.
\]
A directed evaluation of the fixed-bandwidth residual, and point rounding of
$c\widehat v$, add no more than the same amount.  Thus the certificates
satisfy
\begin{equation}
 \label{eq:local-solve-bounds}
 \widehat r_\ell\leq\frac{\zeta}{2^{13}d},
 \qquad
 \widehat s_\ell\leq\frac{\zeta}{2^{14}d},
 \qquad
 \eta_\ell\leq\frac{\zeta}{4096d}.
\end{equation}
The same directed error summation in the normalization step gives
$R_0\leq K_wn(u_p+\tau_{p,E_{\mathrm{sol}}})\leq\zeta/2048$.

We prove the radius certificate by induction.  At stage $\ell$, the exact
identity
\[
 y_\ell-Gy_{\ell-1}
 =(y_\ell-cv_\ell)+cB^{-1}(Bv_\ell-y_{\ell-1})
\]
and $\lVert cB^{-1}\rVert_2\leq1+\zeta/(256d)\leq g_+$ give
$\lVert y_\ell-Gy_{\ell-1}\rVert_2\leq\eta_\ell$.  Therefore
$\lVert y_\ell-G^\ell x\rVert_2\leq R_\ell$ whenever the preceding
certificate is valid.  Moreover,
\[
 \lVert y_{\ell-1}\rVert_2
 \leq\lVert G^{\ell-1}x\rVert_2+R_{\ell-1}
 \leq e^{1/256}+\zeta/128<2,
\]
so the solve estimate used above remains valid at every stage.  Directed
rounding gives $g_+\leq1+\zeta/(256d)+4u_p$, while
\Cref{eq:primitive-precision} gives $du_p<2^{-20}\zeta$.  Hence
\[
 g_+^d<2,
 \qquad
 R_d\leq2(1+8u_p)^d
 \left(\frac{\zeta}{2048}
 +d\frac{\zeta}{4096d}
 +d\tau_{p,E_{\mathrm{sol}}}\right)
 <\frac{\zeta}{128}.
\]
Thus no radius guard fires on a promised input.

For completeness, the rounded factorization represents a matrix within
$\Delta/12$ of $B$, and its accumulated rotations are within $2^{-12}$ of
an isometry.  Its triangular factor therefore has smallest singular value at
least $\Delta/5$, so no back-substitution diagonal vanishes.  A crude reverse
row recurrence places every intermediate below
$(K_wW/\Delta)^{O(n)}$; its exponent has magnitude $O_w(np)$ and fits
$E_{\mathrm{sol}}$.  There are $O_w(n)$ point operations in the
factorization and each solve, and each costs $O(p^2)$ bit operations.  With
$d$ solves, this is $O_w(ndp^2)$ on every branch.
\end{proof}

The parameter bit bounds now show that one common polylogarithmic precision
supports both numerical primitives.

\begin{corollary}[Polylogarithmic guard precision]
\label{cor:primitive-precision}
For the global parameters above, a precision
$p=O_w(L)$ satisfies the precision condition in
\Cref{eq:primitive-precision}.  Successful block quantities have exponent
magnitude $O_w(p)$, successful shifted-solve quantities have exponent
magnitude $O_w(np)$, and the signed exponents themselves use
$O(\log n+\log p)$ bits.
\end{corollary}

\begin{proof}
Each reciprocal on the right side of \Cref{eq:primitive-precision} has
$O(\log(n/(\epsilon\delta)))$ bits by \Cref{eq:parameter-bits}, while the
exact scalar inputs have at most $O(L)$ bits.  Choosing $p$ larger than these
quantities by the fixed factor implicit in $K_w$ proves the precision
assertion.  The exponent bounds were established
in \Cref{lem:banded-bisection,lem:shifted-inverse-power}; encoding a signed
integer exponent of magnitude $O_w(np)$ takes
$O(\log n+\log p)$ bits.
\end{proof}

The localized eigenvalues now determine well-separated shifts, and the
certified inverse-power routine implements their spectral slices.  We define
the exact filters before analyzing their numerical application.

\begin{definition}[Shifted-inverse filter bank]
\label{def:filter-bank}
Let $H$ have eigenpairs $(\lambda_i,u_i)$ in nonincreasing eigenvalue order,
with every two eigenvalues separated by at least $3\Delta$.  Suppose
$|\widetilde\lambda_i-\lambda_i|\leq a$ for $1\leq i\leq k$.  Define
\begin{equation}
 \label{eq:filter-bank-definition}
 z_i=\widetilde\lambda_i+\frac\Delta2,
 \qquad
 F_i=\left[-\frac\Delta2(H-z_iI)^{-1}\right]^d,
 \qquad
 \gamma=\widetilde\lambda_k-\Delta,
 \qquad
 S_\gamma=\sum_{i=1}^kF_i.
\end{equation}
Thus each $F_i$ is a degree-$d$ polynomial in its own shifted inverse; the
sum is a multishift construction for $H$.
\end{definition}

\begin{lemma}[Spectral-step approximation and random eigenvector recovery]
\label{lem:filter-bank}
Under \Cref{def:filter-bank}, exactly $k$ eigenvalues of $H$ are at least
$\gamma$, and
\[
 \left\lVert S_\gamma-\mathbf 1_{[\gamma,\infty)}(H)\right\rVert_2
 \leq\zeta\leq n^{-10}.
\]
Let $M$ be the least power of two with $M\geq16n/\delta$, sample independent
$r_j$ uniformly from $\{-M,-M+1,\ldots,M-1\}$, and let
$x=r/\lVert r\rVert_2$ when $r\ne0$, with $x=e_1$ when $r=0$.
Independently of $H$, with probability at least $1-\delta/16$,
\[
 |u_i^Tx|\geq\beta\qquad(1\leq i\leq n).
\]
On this overlap event, if
$\lVert\widehat y_i-F_ix\rVert_2\leq\zeta/128$ and
$q_i=\widehat y_i/\lVert\widehat y_i\rVert_2$, then, after a possible sign
change,
\begin{equation}
 \label{eq:eigenvector-recovery}
 \lVert q_i-u_i\rVert_2\leq\epsilon/8
 \qquad(1\leq i\leq k).
\end{equation}
\end{lemma}

\begin{proof}
The gap and $a<\Delta/100$ imply
\[
 \lambda_k-\gamma\geq\Delta-a>0,
 \qquad
 \gamma-\lambda_{k+1}\geq2\Delta-a>0
\]
when $k<n$; only the first inequality is needed for $k=n$.  Thus the
threshold has the claimed rank.

By \Cref{eq:spectral-calculus}, the multiplier of $F_i$ at $\lambda_j$ is
\begin{equation}
 \label{eq:filter-multiplier}
 f_i(\lambda_j)=
 \left(\frac{-\Delta/2}{\lambda_j-z_i}\right)^d.
\end{equation}
For $j=i$, write
$\varepsilon_i^{\mathrm{loc}}=\lambda_i-\widetilde\lambda_i$.  Then
\[
 f_i(\lambda_i)=
 \left(1-\frac{2\varepsilon_i^{\mathrm{loc}}}{\Delta}\right)^{-d}.
\]
Since $|\varepsilon_i^{\mathrm{loc}}|\leq a
\leq\zeta\Delta/(2^{10}d)$, the inequality
$|(1-s)^{-d}-1|\leq2d|s|$ for $d|s|\leq1/4$ gives
\begin{equation}
 \label{eq:on-target-filter}
 |f_i(\lambda_i)-1|\leq\zeta/128.
\end{equation}
If $j\ne i$, the $3\Delta$ gap gives
\[
 |\lambda_j-z_i|\geq3\Delta-\Delta/2-a>2\Delta,
 \qquad
 |f_i(\lambda_j)|\leq4^{-d}\leq\frac{\zeta}{128n}.
\]
The one-step operator
\[
 G_i=-\frac\Delta2(H-z_iI)^{-1}
\]
has its target multiplier equal to
$(1-2\varepsilon_i^{\mathrm{loc}}/\Delta)^{-1}$ and every off-target
multiplier below $1/4$ in absolute value.  Hence
\begin{equation}
 \label{eq:filter-power-bound}
 \lVert G_i\rVert_2\leq1+\frac{\zeta}{256d}.
\end{equation}
\Cref{eq:filter-degree,eq:on-target-filter}, summed
over the slices, prove the spectral-step estimate.  They also give the
stronger individual bound
\begin{equation}
 \label{eq:rank-one-filter}
 \lVert F_i-u_iu_i^T\rVert_2\leq\zeta/128.
\end{equation}

For the overlap assertion, fix a unit vector $u$ and choose $j$ with
$|u_j|\geq n^{-1/2}$.  Conditional on all $r_\ell$ with $\ell\ne j$, the
number of values of $r_j$ for which
$|u^Tr|<\beta M\sqrt n$ is at most
$2\beta M\sqrt n/|u_j|+1$.  Consequently
\[
 \Pr\{|u^Tr|<\beta M\sqrt n\}
 \leq\beta n+\frac1{2M}
 \leq\frac\delta{16n}.
\]
Since $\lVert r\rVert_2\leq M\sqrt n$, the complementary event gives
$|u^Tx|\geq\beta$.  A union bound over the eigenbasis proves the simultaneous
claim and excludes $r=0$.

On that event put $\alpha_i=u_i^Tx$.  By
\Cref{eq:rank-one-filter},
$\lVert F_ix-\alpha_iu_i\rVert_2\leq\zeta/128$.  Including the numerical
application error makes the right side at most $\zeta/64$.  Since
$|\alpha_i|\geq\beta$ and $\zeta\leq\epsilon\beta/64$,
\Cref{eq:normalization-stability} bounds the change in direction by
$2(\zeta/64)/\beta\leq\epsilon/16$, which is stronger than
\Cref{eq:eigenvector-recovery}.
\end{proof}

The slices are already nearly rank one, so normalization replaces a dense
range orthogonalization.  It remains to assemble the perturbation,
localization, filter application, and certified Rayleigh values into one
total algorithm.

Let $s$ be the least integer for which
\begin{equation}
 \label{eq:perturbation-grid}
 h=\rho2^{-s}\leq\Delta/8.
\end{equation}
For the input matrix let $b_A$ be the largest binary encoding length of an
entry, and choose the literal integer precision
\begin{equation}
 \label{eq:chosen-precision}
 p=\max\left\{
 \begin{array}{l}
  8,\ b_A+8,\ b_\epsilon+8,\ b_\delta+8,\ s+8,\\[2pt]
  \left\lceil1+\log_2\dfrac{2^{20}K_wW^2}{\nu^2}\right\rceil,\quad
  \left\lceil1+\log_2
   \dfrac{2^{20}K_wn^2dW}{\zeta\Delta}\right\rceil,\\[6pt]
  \left\lceil1+\log_2\dfrac{2^{20}K_wn}{\epsilon}\right\rceil,\quad
  \left\lceil1+\log_2(2^{20}K_wn^2)\right\rceil
 \end{array}
 \right\}.
\end{equation}
The last displayed term makes the third inequality in
\Cref{eq:primitive-precision} explicit; it does not alter the
$O_w(\log(n/(\epsilon\delta)))$ precision bound.

The perturbation law is the one from \Cref{lem:random-gap}: sample independent
uniform $J_j\in\{0,\ldots,2^s-1\}$, set
\begin{equation}
 \label{eq:finite-perturbation}
 D=\operatorname{diag}(hJ_1,\ldots,hJ_n),
 \qquad
 H=A+D.
\end{equation}
The random vector law is the one in \Cref{lem:filter-bank}.  Define the
finite fallback record
\begin{equation}
 \label{eq:fallback-record}
 \mathcal R_k=((0,e_1,e_1),\ldots,(0,e_k,e_k)).
\end{equation}

We also fix the Rayleigh output rule outside the algorithm.  Given a nonzero
dyadic vector $y_i$, rescale it by a common exact power of two until its
largest coordinate has magnitude in $[1/2,1)$.  Under the exponent guard
$E_{\mathrm{ray}}=2E_{\mathrm{sol}}+2p$, let directed evaluation produce
intervals
\begin{equation}
 \label{eq:rayleigh-intervals}
 \mathcal N_i\ni y_i^TAy_i,
 \qquad
 \mathcal D_i\ni y_i^Ty_i,
 \qquad
 I_i=\mathcal N_i/\mathcal D_i.
\end{equation}
When $\inf\mathcal D_i>0$ and
$\operatorname{width}(I_i)\leq\epsilon W/128$, define
\begin{equation}
 \label{eq:rayleigh-output}
 \widetilde\sigma_i=\frac{\inf I_i+\sup I_i}{2},
 \qquad
 \widehat\sigma_i=\max\{0,\widetilde\sigma_i\}.
\end{equation}
Call $(A,k,\epsilon,\delta)$ an \emph{admissible eigenpair input} if $A$ is
real symmetric positive definite of fixed half-bandwidth $w$, its nonzero
entries are integers of $O(\log n)$ bits, its spectrum is contained in
$[n^{-10},n^{10}]$, $1\leq k\leq n$, and
$0<\epsilon,\delta\leq1/4$ are finite dyadics.

\begin{algbox}
\textbf{\uline{\BandedPSDEigenpairs}}: Top eigenpairs of a
fixed-bandwidth positive definite matrix.
\phantomsection\label{alg:main}

\uline{Input}: An admissible tuple $(A,k,\epsilon,\delta)$ as defined above.

\uline{Output}: The $k$ triples $(\widehat\sigma_i,q_i,q_i)$, with
$\widehat\sigma_i$ nonnegative dyadic and $q_i$ a normalized dyadic-vector
encoding.

\uline{Guarantee}: With probability at least $1-9\delta/64$, every bound in
\Cref{thm:main} holds.

\uline{Totality}: Every branch terminates, and any detected failure returns
$\mathcal R_k$ from \Cref{eq:fallback-record}.

\begin{enumerate}
 \item If $n=1$, return $((A_{11},e_1,e_1))$.  Otherwise compute
 $W,b_A,\rho,\Delta,\beta,\zeta,d,a,m_{\mathrm{bis}},N_{\mathrm{qry}},
 \nu,s,h$, and $p$ from
 \Cref{eq:rho-range,eq:chosen-precision}.
 \item Sample $D$ and form the named perturbed matrix $H$ according to
 \Cref{eq:finite-perturbation}.
 \item First for $i=k$, then for $i=1,\ldots,k-1$, invoke
 $\RandomizedBandedBisect(T=H,i=i,W=W,a=a,\nu=\nu,
 m_{\mathrm{bis}}=m_{\mathrm{bis}},p=p)$ with fresh randomness, obtaining
 $(s_i,\widetilde\lambda_i)$.  If $s_i=\mathrm{FAIL}$, return
 $\mathcal R_k$; after the first call set $\gamma=\widetilde\lambda_k-\Delta$.
 \item Compute $M$, sample $r$, and form its normalized encoding $x$ by the
 law in \Cref{lem:filter-bank}.  For each $i=1,\ldots,k$, set
 $z_i=\widetilde\lambda_i+\Delta/2$ and invoke
 $\ShiftedInversePower(T=H,W=W,z=z_i,c=-\Delta/2,\Delta=\Delta,
 \zeta=\zeta,d=d,x=x,p=p)$,
 obtaining $(s_i^{\mathrm{inv}},y_i,R_i^{\mathrm{filt}})$.  If
 $s_i^{\mathrm{inv}}=\mathrm{FAIL}$,
 $R_i^{\mathrm{filt}}>\zeta/128$, or $y_i=0$, return $\mathcal R_k$;
 otherwise store $y_i$ as the normalized encoding of $q_i$.
 \item For every $i$, evaluate the intervals in
 \Cref{eq:rayleigh-intervals}.  If evaluation fails,
 $\inf\mathcal D_i\leq0$, or
 $\operatorname{width}(I_i)>\epsilon W/128$, return $\mathcal R_k$.
 Otherwise compute $\widehat\sigma_i$ by \Cref{eq:rayleigh-output} and
 return the $k$ triples.
\end{enumerate}
\end{algbox}

\begin{lemma}[Correctness and bit complexity of the banded eigensolver]
\label{lem:main-algorithm}
\hyperref[alg:main]{\BandedPSDEigenpairs} has the output guarantee,
totality, worst-case $O_w(nkL^3)$ bit complexity,
$O(nL+kL^2)$ random-bit usage, and $O(nkL)$ output size stated in
\Cref{thm:main}.
\end{lemma}

\begin{proof}
Let $\mathcal E_{\mathrm{gap}}$ be the gap event in
\Cref{lem:random-gap}, let $\mathcal E_{\mathrm{bis}}$ be the event that all
adaptive bisection queries are $2\nu$-safe, and let
$\mathcal E_{\mathrm{ov}}$ be the overlap event in
\Cref{lem:filter-bank}.  By \Cref{lem:random-gap,lem:banded-bisection},
\begin{equation}
 \label{eq:filter-event-probability}
 \Pr(\mathcal E_{\mathrm{gap}}\cap\mathcal E_{\mathrm{bis}})
 \geq1-\frac\delta{16}-\frac\delta{64}
 =1-\frac{5\delta}{64}.
\end{equation}
The overlap vector is independent, and a union bound gives
\begin{equation}
 \label{eq:main-event-probability}
 \Pr(\mathcal E)\geq1-\frac{9\delta}{64}>1-\delta,
 \qquad
 \mathcal E=\mathcal E_{\mathrm{gap}}\cap
 \mathcal E_{\mathrm{bis}}\cap\mathcal E_{\mathrm{ov}}.
\end{equation}

The choice in \Cref{eq:chosen-precision} implies the sufficient condition in
\Cref{eq:primitive-precision}.  All exact encodings of $D$, $H$, the scalar
parameters, bisection lattice endpoints, and shifts have at most
\[
 b_A+s+8+\log_2(W/\Delta)+\log_2(1/\zeta)
 +\log_2(W/a)+\log_2(W/\nu)
\]
bits.  By \Cref{eq:parameter-bits} and the explicit terms in
\Cref{eq:chosen-precision}, this is at most $8p$, hence at most
$K_wp/32$.  Their exponent magnitudes have the same margin below
$E_{\mathrm{blk}}/16$ and $E_{\mathrm{sol}}/16$.  The representative $r$
and the integer $d$ satisfy these bounds as well.  Thus every primitive call
meets its literal input premises.

On $\mathcal E$, every bisection returns an $a$-accurate
$\widetilde\lambda_i$.  By \Cref{lem:filter-bank}, the interval
$[\gamma,2W]$ contains exactly the top $k$ eigenvalues of $H$ and the exact
filter sum satisfies
\begin{equation}
 \label{eq:algorithm-step-filter}
 \left\lVert S_\gamma-
 \mathbf 1_{[\gamma,\infty)}(H)\right\rVert_2
 \leq\zeta\leq n^{-10}.
\end{equation}
Together with \Cref{eq:filter-event-probability}, this also proves the
conditional projector guarantee independently of the overlap event.  For every slice,
\[
 \operatorname{dist}(z_i,\operatorname{spec}(H))
 \geq\Delta/2-a>\Delta/3,
\]
while $|z_i|\leq3W$, $|c|=\Delta/2$, and the power promise is precisely
\Cref{eq:filter-power-bound}.  Therefore
\Cref{lem:shifted-inverse-power} gives a numerical application error at most
$\zeta/128$.  The recovery part of \Cref{lem:filter-bank} then gives, after
signs are chosen,
\begin{equation}
 \label{eq:q-close-u}
 \lVert q_i-u_i(H)\rVert_2\leq\epsilon/8.
\end{equation}
In particular, no $y_i$ is zero: before numerical error,
\Cref{eq:rank-one-filter} and $|u_i^Tx|\geq\beta$ give
$\lVert F_ix\rVert_2\geq\beta-\zeta/128$, and the application error is
less than $\beta/4$.

Put $\theta_i=\lambda_i(H)$.  Since
$\lVert H\rVert_2\leq W+\rho<2W$, the spectral theorem and
\Cref{eq:q-close-u} give
\begin{equation}
 \label{eq:H-residual}
 \lVert Hq_i-\theta_iq_i\rVert_2
 \leq2\lVert H\rVert_2\lVert q_i-u_i(H)\rVert_2
 \leq\epsilon W/2.
\end{equation}
Because $A=H-D$ and $\lVert D\rVert_2\leq\rho$,
\begin{equation}
 \label{eq:pre-rayleigh-residual}
 \lVert Aq_i-\theta_iq_i\rVert_2
 \leq\epsilon W/2+\rho<\epsilon W.
\end{equation}

We next show that every Rayleigh certificate succeeds on $\mathcal E$.
After its exact common rescaling, $D_i=y_i^Ty_i$ satisfies
$1/4\leq D_i\leq n$.  The row-sum definition of $W$ and
$2|y_{i,r}y_{i,s}|\leq y_{i,r}^2+y_{i,s}^2$ imply
\[
 \sum_{r,s}|A_{rs}y_{i,r}y_{i,s}|\leq WD_i.
\]
Sequential directed dot products under $E_{\mathrm{ray}}$ enclose the
numerator and denominator with endpoint errors at most
\[
 8K_wn(u_p+\tau_{p,E_{\mathrm{ray}}})WD_i
 \quad\text{and}\quad
 8K_wn(u_p+\tau_{p,E_{\mathrm{ray}}})D_i,
\]
respectively.  Since $\tau_{p,E_{\mathrm{ray}}}\leq u_p$, the Rayleigh term
in \Cref{eq:chosen-precision} makes each relative error at most
$\epsilon/2^{15}$.  Thus $\inf\mathcal D_i>D_i/2\geq1/8$, and directed
division gives $\operatorname{width}(I_i)\leq\epsilon W/128$.  Hence no
Rayleigh fallback occurs and
\begin{equation}
 \label{eq:rayleigh-midpoint-error}
 |\widetilde\sigma_i-q_i^TAq_i|\leq\epsilon W/256.
\end{equation}

For the exact Rayleigh quotient $\sigma_i=q_i^TAq_i$, write
$r_i=Aq_i-\theta_iq_i$.  Then
\[
 Aq_i-\sigma_iq_i=r_i-(q_i^Tr_i)q_i,
\]
so \Cref{eq:pre-rayleigh-residual} gives
$\lVert Aq_i-\sigma_iq_i\rVert_2<2\epsilon W$.  Since $A\succeq0$,
$\sigma_i\geq0$; clamping the approximate Rayleigh value to zero cannot
increase its error.  By \Cref{eq:rayleigh-midpoint-error} and
\Cref{eq:W-comparison},
\begin{equation}
 \label{eq:final-residual}
 \lVert Aq_i-\widehat\sigma_iq_i\rVert_2
 \leq C_w\epsilon\lVert A\rVert_2.
\end{equation}

The eigenvalue estimate follows explicitly from
\begin{align*}
 |\widehat\sigma_i-\lambda_i(A)|
 &\leq |\widehat\sigma_i-q_i^TAq_i|
   +|q_i^TAq_i-q_i^THq_i|\\
 &\quad+|q_i^THq_i-\lambda_i(H)|
   +|\lambda_i(H)-\lambda_i(A)|.
\end{align*}
The four terms are at most $\epsilon W/256$, $\rho$,
$2\lVert H\rVert_2\lVert q_i-u_i(H)\rVert_2$, and $\rho$, respectively;
the last bound follows from Weyl's inequality in \Cref{eq:weyl}.
\Cref{eq:rho-range,eq:q-close-u,eq:W-comparison} therefore yield
\begin{equation}
 \label{eq:final-eigenvalue}
 |\widehat\sigma_i-\lambda_i(A)|
 \leq C_w\epsilon\lVert A\rVert_2.
\end{equation}

For $i\ne j$, the simple-spectrum eigenvectors of $H$ are orthogonal, and
\Cref{eq:q-close-u} gives
\[
 |q_i^Tq_j|
 \leq\lVert q_i-u_i(H)\rVert_2+
 \lVert q_j-u_j(H)\rVert_2+
 \lVert q_i-u_i(H)\rVert_2\lVert q_j-u_j(H)\rVert_2
 \leq C\epsilon.
\]
The encoded columns are exactly unit vectors.  Hence the maximum absolute
row sum of $Q^TQ-I$ is at most $C(k-1)\epsilon$, and symmetry gives
\begin{equation}
 \label{eq:gram-bound}
 \lVert Q^TQ-I\rVert_2\leq Ck\epsilon\leq Cn\epsilon.
\end{equation}

It remains to count bits.  Computing the scalar parameters costs
$O(nL+L^2)$ bit operations.  The input assumption gives
$b_A=O(\log n)$, while \Cref{eq:parameter-bits} and
\Cref{eq:chosen-precision} give $s=O(L)$ and $p=O_w(L)$.  Each of the $k$
bisections has $O(L)$ iterations, and each of the $k$ filters has $d=O(L)$
shifted solves.  By
\Cref{lem:banded-bisection,lem:shifted-inverse-power}, these operations cost
\[
 O_w(nkLp^2)=O_w(nkL^3)
\]
bit operations.  The primitive bounds count every directed endpoint,
division, square root, exponent-guard update, and residual-certificate
operation among these $O_w(nkL)$ guarded scalar operations; each such
operation costs $O(p^2)$ bits.  Scalar bookkeeping and the $k$ certified
Rayleigh quotients are smaller.
Perturbation and overlap sampling use $O(nL)$ random
bits, while \Cref{lem:banded-bisection} gives $O(L^2)$ random bits for each
of the $k$ localizations.  Thus the total is $O(nL+kL^2)$ random bits.  All
iteration counts are explicit and every exceptional arithmetic operation
returns the finite fallback, so the bit-operation bound holds on every random
branch.

Every returned coordinate has a $p$-bit significand and an exponent of
magnitude $O_w(np)$.  Encoding the signed exponent uses
$O(\log n+\log p)=O(L)$ bits.  Exact common rescaling, directed endpoints,
and the dyadic midpoint preserve this bound, so the $k$ normalized-vector
records and eigenvalue records use $O(nkL)$ bits.  Combining this accounting
with \Cref{eq:main-event-probability} proves every claim.
\end{proof}

\Cref{lem:main-algorithm} proves \Cref{thm:main}.

\bibliographystyle{alpha}
\bibliography{references}

\end{document}
